\documentclass[11pt]{article}
\usepackage[a4paper,margin=25mm]{geometry}
\usepackage{amsmath,amssymb,amsthm,mathtools,booktabs,hyperref}
\setlength{\emergencystretch}{2em}
\newtheorem{theorem}{Theorem}
\newtheorem{lemma}[theorem]{Lemma}
\newtheorem{proposition}[theorem]{Proposition}
\newtheorem{corollary}[theorem]{Corollary}
\newcommand{\Q}{\mathbb Q}\newcommand{\Z}{\mathbb Z}
\newcommand{\N}{\mathbb Z_{\ge0}}\newcommand{\A}{\mathcal A}
\newcommand{\trdeg}{\operatorname{trdeg}}
\title{Absorbing the mixed-integer wrapper into a finite algebra\\
Rational undecidability in 82, 100, and 52 variables}
\author{\(\mathbb Z\mathbb Q\) Diophantine Project}
\date{23 November 2026}
\begin{document}
\maketitle

\begin{abstract}
The mixed-integer universal equation uses eight integer inputs and one
nonnegative integer input.  The previous construction imposed nine
integer predicates and represented the last input by three rational
coordinates.  We retain all nine integer predicates, but absorb the
universal equation and one of the three coordinates into the joint
finite algebra.  The added algebra has rank sixteen and is generated
by one element.  This reduces the common-numerator constructions from
83, 101, and 53 coordinates to 82, 100, and 52, respectively.  The last
number retains ordinary rational \(abc\); the first two introduce no
\(abc\) assumption.  All three retain the arithmetic premises of their
underlying integer predicates.  We give the exact polynomial recipe,
both directions of the equivalence, denominator bookkeeping, and
certified total-degree bounds \(2^{650}\), \(2^{714}\), and \(2^{500}\).
The final polynomials have unrestricted degree; no quartic conclusion
is drawn.
\end{abstract}

\section{The source theorem and the improvement}

Bayer, David, Ha{\ss}ler, Matiyasevich and Schleicher
\cite[Definition 5.3 and Theorem 3]{BDHMS} give an explicit construction
with the following property.  Given
\(P(a,z_1,\ldots,z_\nu)\in\Z[a,z_1,\ldots,z_\nu]\), it constructs
\[
 Q(a,f,g,h,k,l,w,x,y,n)\in\Z[a,f,g,h,k,l,w,x,y,n]
\]
such that, for every \(a\in\N\),
\begin{equation}\label{source}
 \exists z\in\N^\nu\ P(a,z)=0
 \quad\Longleftrightarrow\quad
 \exists f,g,h,k,l,w,x,y\in\Z,\ n\in\N\ Q(a,f,g,h,k,l,w,x,y,n)=0.
\end{equation}
This refines the mixed-domain construction of
Sun \cite[Theorem 1.1(i)]{Sun}.  The integer predicates used in the
project have a finite-root presentation: ordinary rational coordinates
specify a finite algebra, and a guarded rational character of that
algebra is equivalent to a successful predicate.

The conjunction of nine such predicates has already been constructed
with one joint finite-algebra encoder and with common reference and
numerator points \cite{Sharing}.  Its ordinary coordinates include the
nine inputs.  The nonnegative input was then replaced by
\[
                 n=s^2+t^2+u^2+u.
\]
The improvement in this paper removes \(n\) from the ordinary coordinates
and retains only \(s,t\) from this wrapper.  Both \(n\) and the missing
\(u\) are recovered from finite algebra generators.  The final
single-coordinate encoder handles these generators together with all
previous finite witnesses.

\begin{theorem}\label{main}
Apply the construction below to the three common-numerator joint
integer predicates.  It gives rational undecidability bounds of
\[
                 82,\qquad 100,\qquad 52
\]
variables, respectively.  The first route retains its sufficiently
large fixed, possibly ineffective parameter; the second uses the
explicit three-detector parameter \(L=143\); the third retains ordinary
rational \(abc\) and its sufficiently large fixed odd parameter.
The respective integer-coefficient polynomials may be chosen with
total degrees strictly less than
\[
                       2^{650},\qquad 2^{714},\qquad 2^{500}.
\]
No integer predicate is omitted.  These conclusions retain exactly the
arithmetic premises of the corresponding joint predicates, and are
independent of the quarantined quartic universality transcription.
\end{theorem}

The extra finite algebra has rank \(16\), not sixteen new quantified
coordinates.  Its rational characters are handled by the existing
single-coordinate encoder.

\section{An exact recipe for the mixed-integer polynomial}

We specify the source to make the crucial leading coefficient
auditable.  The notation follows \cite[Definitions 2.1, 4.1 and 5.3]{BDHMS},
with subscripts to distinguish different auxiliary expressions.

First replace \(P\) by
\[
 \overline P(a,z_1,\ldots,z_{\nu+1})
       =P(a,z_1,\ldots,z_\nu)^2+(z_{\nu+1}-1)^2.
\]
Put \(v=\nu+1\), let \(\Delta=\deg\overline P\), and write
\(\overline P=\sum_I c_I Z^I\), where \(I=(I_0,\ldots,I_v)\).
Set
\[
 L_P=\sum_I|c_I|,\qquad n_j=(\Delta+1)^j,
\]
and choose the smallest power \(\beta=4^r\) satisfying
\[
                       \beta>2\Delta!L_P(v+2)^\Delta.
\]
Define
\[
 \alpha=\Delta(\Delta+1)^v+1,\qquad
 \gamma=\beta^{(\Delta+1)^v},\qquad
 b=1+3(2a+1)f,\qquad \mu=\gamma b^\alpha,\qquad B_0=\beta b^\Delta.
\]
For the mask set
\[
 M_0=\sum_{j=0}^{n_v}m_j B_0^j,\qquad
 m_j=
 \begin{cases}
 B_0-b,&j\in\{n_1,\ldots,n_v\},\\
 B_0-1,&\text{otherwise}.
 \end{cases}
\]
The coefficient-encoding polynomial and the value polynomial are
\[
 \operatorname{coeffs}(B_0)
 =\sum_{\lVert I\rVert\le\Delta}
 I!(\Delta-\lVert I\rVert)!\,c_I\,
 B_0^{\,n_{v+1}-\sum_{j=0}^{v}I_jn_j},
\]
\[
 c=1+aB_0+g,\qquad
 K_0=c^\Delta\operatorname{coeffs}(B_0)
           +\frac{B_0}{2}\sum_{j=0}^{(2\Delta+1)n_v}B_0^j.
\]
The displayed exponents are nonnegative.  Since \(\beta\) is a power
of four, \(B_0/2\) is an integer-coefficient polynomial.  Finally put
\[
 \begin{gathered}
 N_0=B_0^{n_v+1},\qquad
 N_1=4B_0^{(2\Delta+1)n_v+1},\qquad N=N_0N_1,\\
 S_0=g+2K_0N_0,\qquad
 T_0=M_0+(B_0-2)B_0^{n_{v+1}}N_0,\\
 R_0=(S_0+T_0+1)N+T_0+1,\qquad
 X=(N-1)R_0,\qquad Y=N^2.
 \end{gathered}
\]
Thus \(b,\mu,X,Y\) are explicitly specified integer-coefficient
polynomials in \(a,f,g\).  The remaining bridge polynomials are
\begin{align*}
 U&=2lXY,& V&=4gwY,\\
 A&=U(V+1),& B_1&=2X+1,\\
 C&=B_1+(A-2)h,& D&=(A^2-4)C^2+4,\\
 E&=C^2Dx,& F&=4(A^2-4)E^2+1,\\
 G&=1+CDF-2(A+2)(A-2)^2E^2,\\
 H&=C+B_1F+(2y-1)CF,\\
 I&=(G^2-1)H^2+1,& J&=X+1+k(U^2V-2).
\end{align*}
Define
\begin{equation}\label{M3data}
 \begin{gathered}
 A_1=b,\qquad A_2=DFI,\qquad A_3=(U^4V^2-4)J^2+4,\\
 S=2A-5,\qquad T=3bwC-2(b^2w^2-1),\\
 R=f^2l^2x^2
       \left[8\mu^3gJ^2
       -g^2\left\{32(C-lJY)^2\mu^3+g^2J^2\right\}\right],\\
 Z_0=1+A_1^2+A_2^2+A_3^2.
 \end{gathered}
\end{equation}
None of these expressions depends on \(n\).  The source polynomial is
\begin{equation}\label{Qproduct}
 Q(n)=
 \prod_{\varepsilon\in\{\pm1\}^3}
 \left[
 S^2n+T^2
 -S^2(2R-1)\left(
 T^2+Z_0^3+
 \varepsilon_1\sqrt{A_1}
 +\varepsilon_2Z_0\sqrt{A_2}
 +\varepsilon_3Z_0^2\sqrt{A_3}
 \right)\right].
\end{equation}
The product is an integer polynomial: changing each formal square
root's sign permutes its factors, so only even powers of that square
root remain.  Formula \eqref{Qproduct} is a coefficient prescription;
it does not introduce three more ordinary variables.

\begin{lemma}\label{leading}
The degree of \(Q\) in \(n\) is exactly \(8\), and its leading
coefficient is \(S^{16}\).  Moreover, \(S\ne0\) at every integral
specialization of \(a,f,g,h,k,l,w,x,y\).
\end{lemma}
\begin{proof}
There are eight factors in \eqref{Qproduct}, each with leading
coefficient \(S^2\) in \(n\).  The polynomial \(S\) is not identically
zero.  At an integral specialization, \(A\) is an integer, hence
\(2A-5\) is odd and cannot vanish.  In fact \(U\), and therefore \(A\),
is even, so \(S\equiv3\pmod4\).
\end{proof}

\section{The monic rank-sixteen wrapper}

Introduce a formal variable \(m\), and define
\begin{equation}\label{Pm}
 \mathcal P(m)=
 \prod_{\varepsilon\in\{\pm1\}^3}
 \left[
 m+T^2
 -S^2(2R-1)\left(
 T^2+Z_0^3+
 \varepsilon_1\sqrt{A_1}
 +\varepsilon_2Z_0\sqrt{A_2}
 +\varepsilon_3Z_0^2\sqrt{A_3}
 \right)\right].
\end{equation}
It is monic of degree eight, with integer-polynomial coefficients,
and satisfies the exact identity
\begin{equation}\label{Pidentity}
                         \mathcal P(S^2n)=Q(n).
\end{equation}

For a prescription involving no formal radicals, put
\[
 \kappa=S^2(2R-1),\qquad z=m+T^2-\kappa(T^2+Z_0^3),
\]
\[
 a_1=\kappa^2A_1,\qquad
 a_2=\kappa^2Z_0^2A_2,\qquad
 a_3=\kappa^2Z_0^4A_3.
\]
Then
\begin{equation}\label{radicalfree}
 \mathcal P(m)=
 \left[(z^2+a_3-a_1-a_2)^2+4z^2a_3-4a_1a_2\right]^2
 -16z^2a_3(z^2+a_3-a_1-a_2)^2.
\end{equation}
Indeed, first take the product over the signs of
\(\sqrt{a_1},\sqrt{a_2}\), which is
\(((z+\sqrt{a_3})^2-a_1-a_2)^2-4a_1a_2\), and then multiply
its two conjugates with respect to \(\sqrt{a_3}\).

Let \(R_{\mathrm{base}}\) be the polynomial ring in the retained
ordinary coordinates, including the eight integer-source coordinates
and two new coordinates \(s,t\).  Adjoin
\begin{equation}\label{wrapperalgebra}
 \mathcal B=
 R_{\mathrm{base}}[m,\xi]/
 \left(\mathcal P(m),\
       \xi^2+S\xi+S^2(s^2+t^2)-m\right).
\end{equation}
Both equations are monic in the newly adjoined generator.  Hence
\(\mathcal B\) is free of rank \(8\cdot2=16\), with basis
\[
                    \{m^j,m^j\xi:0\le j<8\}.
\]
Also
\[
                     m=\xi^2+S\xi+S^2(s^2+t^2),
\]
so \(\xi\) alone generates this algebra over its base.  Equivalently,
it is the monic degree-sixteen algebra
\[
 R_{\mathrm{base}}[\xi]/
 \left(\mathcal P\bigl(\xi^2+S\xi+S^2(s^2+t^2)\bigr)\right).
\]
The two-stage basis is more convenient for degree estimates.

Add \(S\) to the selected-character guard.  At any rational character
with \(S\ne0\), define
\begin{equation}\label{recovery}
                         n=\frac m{S^2},\qquad u=\frac{\xi}{S}.
\end{equation}
Then \eqref{wrapperalgebra} gives
\begin{equation}\label{recovered}
          Q(n)=0,\qquad n=s^2+t^2+u^2+u.
\end{equation}

\begin{lemma}[The integer nonnegative wrapper]\label{nonnegative}
For an integer \(n\),
\[
 n\ge0
 \quad\Longleftrightarrow\quad
 n=s^2+t^2+u^2+u
 \text{ for some }s,t,u\in\Q.
\]
In the forward direction the three witnesses can be integers.
\end{lemma}
\begin{proof}
The displayed expression is
\(s^2+t^2+(u+\tfrac12)^2-\tfrac14\), so it is at least
\(-\tfrac14\).  An integral value is therefore nonnegative.
Conversely, for \(n\ge0\), the integer \(4n+1\), being congruent to
\(1\) modulo four, is not of the excluded form \(4^a(8b+7)\) in the
three-squares theorem.  Write it as three integer squares.  Exactly
one of the three integers is odd; write it as \(2u+1\), and the two
even integers as \(2s,2t\).  Division by four gives the identity.
This is also \cite[Lemma 6.1]{BDHMS}.
\end{proof}

Crucially, this lemma is applied only after the ninth integer predicate
has certified \(n=m/S^2\in\Z\).  A rational root of \(\mathcal P\)
does not, by itself, imply that \(m/S^2\) is integral.  The original
divisibility information encoded by \(M_3\) is not being discarded.

\section{Merging with the nine integer predicates}

The joint common-numerator sources have ordinary input coordinates
\(f,g,h,k,l,w,x,y,n\), ordinary elliptic-chart and pole parameters,
point blocks of rank four, and terminal monic quadratic pole blocks.
Their retained guards contain no terminal pole roots.  Each point
block has basis \(1,U_0,V_0,W_0\), with coefficient polynomials
\(A_0,B_0\) depending only on that point's retained coordinates.
In particular those coefficients are independent of all nine inputs.
This last assertion is checked directly on every actual source DAG
by the accompanying verifier.

Start with the wrapper algebra \eqref{wrapperalgebra}.  Adjoin the old
point blocks once each, substitute \(n=m/S^2\) in every other coefficient
and guard, and then adjoin the terminal quadratic roots, clearing
denominators monically.  The order is important: a pole coefficient
may now depend on \(m\), but it depends on no subsequent pole root.

For example, suppose an old pole satisfies
\[
                              q^2+c(n)=0
\]
and \(e\) is an upper bound for the degree of \(c\) in \(n\).
Then
\[
                       \widetilde c=S^{2e}c(m/S^2)
\]
is a polynomial element of the preceding algebra.  Replacing the finite
root by \(\widetilde q=S^e q\) gives the monic equation
\[
                           \widetilde q^{\,2}+\widetilde c=0.
\]
It is reversible at every selected character since \(S\ne0\).
For each old polynomial guard \(g_j(n)\), choose an upper bound \(e_j\)
for its degree in \(n\) and retain
\[
                    \widetilde g_j=S^{2e_j}g_j(m/S^2).
\]
The single new guard is
\[
                                  S\prod_j\widetilde g_j.
\]
It selects exactly the old allowed characters with \(S\ne0\).
No denominator is required to be nonzero at every unselected character.
Zero terminal pole roots remain permitted.

\begin{proposition}\label{equivalence}
At a fixed natural-number parameter \(a\), a guarded rational character
of the enlarged finite algebra exists for some ordinary retained
coordinates if and only if the right side of \eqref{source} holds.
\end{proposition}
\begin{proof}
Suppose first that the right side of \eqref{source} holds.  The nine
integer inputs admit a simultaneous witness in the appropriate
common-numerator joint predicate.  By Lemma~\ref{leading}, \(S\ne0\).
Choose integers \(s,t,u\) from Lemma~\ref{nonnegative}, and put
\(m=S^2n\), \(\xi=Su\).  Identity \eqref{Pidentity} and the wrapper
identity show that these satisfy \eqref{wrapperalgebra}.  Rescale every
old terminal root by its prescribed power of \(S\).  Every new guard
is nonzero at this character, because its old guard and \(S\) are
nonzero.  Thus the enlarged algebra has a permitted rational character.

Conversely, a permitted rational character gives \(S\ne0\), so all
rescalings can be inverted.  The old joint predicate then certifies
each of \(f,g,h,k,l,w,x,y,n=m/S^2\) as an integer.  The wrapper gives
\(Q(n)=0\) and the identity in \eqref{recovered}.  By
Lemma~\ref{nonnegative}, \(n\ge0\).  Thus all hypotheses on the right
side of \eqref{source} are satisfied.  Applying the published
mixed-integer theorem yields the left side of \eqref{source}.
\end{proof}

The one-coordinate guarded-character theorem \cite{Encoder} applies
to the enlarged finite free algebra.  It produces one
integer-coefficient polynomial.  There is no remaining external source
equation or external positivity equation to combine by another sum of
squares.

\section{The coordinate ledger}

The nine integer predicates still have nine input arguments, but their
last argument \(m/S^2\) is now an element of the finite algebra.  It is
not an ordinary quantified coordinate.  The previous common numerator
point is retained, and the private root point for each of the nine
inputs is also retained.

\begin{center}
\begin{tabular}{lrrr}
\toprule
Ordinary rational coordinates & General & Explicit & Ordinary \(abc\)\\
\midrule
Eight unrestricted integer-source inputs &8&8&8\\
Shared elliptic-point coordinates &16&16&4\\
Private root-point coordinates &18&18&18\\
Common contact coordinate &1&1&1\\
Private pole parameters &36&54&18\\
Two remaining three-squares parameters &2&2&2\\
One joint finite-algebra encoder &1&1&1\\
\midrule
Total &82&100&52\\
\bottomrule
\end{tabular}
\end{center}

The finite roots \(m,\xi\), the point-algebra generators, and the
terminal pole roots are internal to the encoder.  They add no entries
to this table.  The original finite-algebra ranks were
\(2^{52},2^{61},2^{31}\); the new ranks are
\[
                              2^{56},\qquad2^{65},\qquad2^{35}.
\]
Since \(\xi\) generates both new wrapper stages, the numbers \(\ell\)
of guarded separating elements increase by one, to \(54,63,33\).

\section{An explicit uniform degree bound}

We recall the elementary bound proved in \cite{Encoder}.  If a finite
free algebra has rank \(r\), if \(\ell\) is the number of elements
\(g,g\theta_1,\ldots,g\theta_{\ell-1}\), and if the entries of their
multiplication matrices have total degree at most \(d\), set
\[
 \mathcal J=\ell(r+1)^2,\qquad
 \mathcal D(r,\ell,d)=
 2r(\mathcal J+1)
 \left[1+d\{1+2\mathcal J(1+\ell(\mathcal J+1))\}\right].
\]
The final single-coordinate polynomial has degree at most
\(\mathcal D(r,\ell,d)\).  Nilpotents and exceptional specializations
do not invalidate this bound.

For the classical universal pair of 58 unknowns and degree
four \cite{Jones}, the construction in \cite[Theorem A and Section 6.1]{BDHMS}
supplies the bound
\begin{equation}\label{delta}
 \begin{split}
 \delta&=15616+233856\cdot4
       +(233952\cdot4+467712\cdot4^2)\,9^{59}\\
 &<2^{211}.
 \end{split}
\end{equation}
If the classical source is stated for positive integers, the shifts
\(z_i\mapsto z_i+1\) put it over \(\N\) without changing its degree
or its number of unknowns.  The displayed value is a bound for the source
after the original ternary substitution,
so it also bounds the degrees of \(Q\) and \(S\).  It bounds the
coefficients of \(\mathcal P\) as well.  Indeed, writing
\[
 \mathcal P(m)=\sum_{j=0}^{8}p_jm^j,\qquad
 Q(n)=\sum_{j=0}^{8}q_jn^j,
\]
identity \eqref{Pidentity} gives
\(q_j=S^{2j}p_j\).  In an integral polynomial ring,
\(\deg p_j\le\deg q_j\le\delta\) whenever \(p_j\ne0\).

Consequently multiplication by \(m\), in the standard rank-eight
basis, has matrix entries of degree at most \(\delta\).
In the rank-sixteen basis used above, multiplication by \(\xi\) has
block matrix
\[
 \begin{pmatrix}
 0&M_m-S^2(s^2+t^2)I_8\\
 I_8&-SI_8
 \end{pmatrix}.
\]
Its entries have degree at most \(2\delta+2\).

Let \(D\) be one of the old coefficient-DAG degree bounds, evaluated
with ordinary inputs of weight one and point generators weighted by
their multiplication-matrix degrees.  It bounds the degree in the last
input \(n\), as well as the degree of every multiplication-matrix
entry.  Since point-block coefficients do not depend on \(n\), the
cleared expression
\[
                              S^{2D}f(m/S^2)
\]
has multiplication-matrix degree at most \(D(2\delta+1)\).
To see this, an entry monomial with \(n\)-exponent \(j\) has remaining
ordinary-coordinate degree at most \(D-j\).  After substitution and
clearing, its degree is at most
\[
 j\delta+2(D-j)\delta+(D-j)\le D(2\delta+1).
\]
The same estimate follows operation by operation through the
coefficient DAG, without expanding the matrices.

Write \(d_g\) for the old guard bound and \(d_\theta\) for the largest
old separating-generator bound.  Rescaling the pole roots and clearing
the guards as above gives
\[
 d'_g\le d_g(2\delta+1)+\delta,\qquad
 d'_\theta\le d_\theta(2\delta+1).
\]
The second bound also dominates \(2\delta+2\), since every route has
\(d_\theta\ge100232\).  Therefore it is enough to use
\begin{equation}\label{newd}
                      d'=(d_g+d_\theta)(2\delta+1)+\delta.
\end{equation}
Only one factor \(S\) is needed in the product guard; all other cleared
denominators are powers of that same nonzero factor.

\begin{center}
\begin{tabular}{lrrrr}
\toprule
Route & \(r\) & \(\ell\) & \(d_g+d_\theta\) & Final degree\\
\midrule
General &\(2^{56}\)&54&2302058&\(<2^{650}\)\\
Explicit &\(2^{65}\)&63&4126214&\(<2^{714}\)\\
Ordinary \(abc\)&\(2^{35}\)&33&3309916&\(<2^{500}\)\\
\bottomrule
\end{tabular}
\end{center}

These are exact integer upper-bound calculations using
\(\mathcal D(r,\ell,d')\).  Unlike the earlier wrapper, this construction
does not multiply the final degree bound by four: the source equation
and the quadratic positivity identity are already relations of the
finite algebra.  The family parameters occur only in fixed
coefficients, so the estimates are uniform in their admissible fixed
values.

\section{A general equation-absorption principle}

The preceding construction is an especially simple instance of a
general algebraic device.

\begin{proposition}[Absorbing a finite equation]\label{absorb}
Suppose a rational predicate has ordinary coordinates \((t,x)\),
a triangular finite-root presentation, and a selected-character
nonzero guard.  Impose in addition
\[
              P(t,x)=a_dx^d+\cdots+a_0=0,\qquad a_d\ne0,\quad d\ge1,
\]
where the coefficients are integer polynomials in \(t\).
Then \(x\) can be removed from the ordinary coordinates by adding a
finite monic stage of rank \(d\), keeping one final joint encoder.
\end{proposition}
\begin{proof}
Adjoin \(\eta=a_dx\) by the monic relation
\[
 \eta^d+\sum_{j=0}^{d-1}a_j a_d^{d-1-j}\eta^j=0,
\]
retain \(a_d\) in the guard, and substitute \(x=\eta/a_d\) throughout
the old presentation.  Clear all new denominators by monic rescaling
in the triangular order.  At a selected character the process is
reversible, and every monic extension remains finite free.  Apply the
single-coordinate guarded-character encoder to the entire algebra.
\end{proof}

One may use the actual nonzero degree on finitely many coefficient
strata, recording the vanishing higher coefficients as polynomial
conditions and taking a finite disjunction.  However, a stratum on
which \(P(t,x)\) is identically zero in \(x\) does not by itself permit
eliminating \(x\).  It requires a separate construction or another
finite equation.  The present application avoids this issue because
\(S\ne0\) on every certified integer-source tuple, and its monic wrapper
has degree sixteen at every base specialization.

The saving comes from merging a finite algebraic constraint with the
existing encoder.  It does not establish that a source integer
variable has become automatically integral, nor does it justify
dropping any of the nine integer predicates.

\section{Reproducibility and status}

The checker
\nolinkurl{undecidability/verify_finite_wrapper.py}
verifies the radical-free norm identity, the monic degrees and leading
coefficients, the exact wrapper identity, the independence of all
point coefficients from the nine inputs in the actual source DAGs,
and the three numerical degree and coordinate ledgers.  Its report is
\nolinkurl{undecidability/certificates/finite_wrapper_transfer.json}.
The coefficient prescription consists of
\eqref{M3data}--\eqref{Qproduct},
\eqref{radicalfree}, \eqref{wrapperalgebra}, the pinned common-numerator
DAGs, the specified rescalings, and the determinant formula of
\cite{Encoder}.  Expanding the final polynomial is unnecessary and
would be impractical.

The executable compiler
\nolinkurl{undecidability/code/build_finite_wrapper.py} accepts sparse
integer coefficients of any supplied \(P(a,z_1,\ldots,z_\nu)\).
It constructs the complete BDHMS bridge, the eight monic coefficients,
and every transformed coefficient and guard of the joint source.
The long geometric sums are evaluated as polynomial circuits by binary
doubling.  For example, the mask is constructed as
\[
 (B_0-1)\sum_{j=0}^{n_v}B_0^j
       -(b-1)\sum_{i=1}^{v}B_0^{n_i},
\]
so its circuit has size polynomial in the bit lengths of the exponents.
Enormous fixed powers remain power nodes.

The separate checker
\nolinkurl{undecidability/verify_finite_wrapper_compiler.py}
exactly replays each homogenization instruction and checks the
triangular finite-stage dependencies, sparse coding constants, ranks,
guards and coordinate counts.  It also performs explicitly labelled
finite-field diagnostics of all coefficient nodes and of the
multiquadratic norm.  These diagnostics supplement, rather than
replace, the exact identities proved above.  The stored example
\(P(a,z_1)=z_1^2-a\) is a compiler fixture, not a universal polynomial
or an undecidability witness.  A 58-variable degree-four sparse
stress fixture verifies that the large coding exponents can be
handled without their expansion.

The argument proves a further reduction for unrestricted degree.
The separate quartic checkpoint still carries its explicit
universality hypothesis; this paper neither imports that hypothesis
nor resolves its source discrepancy.

\begin{thebibliography}{9}
\bibitem{BDHMS}
J. Bayer, M. David, M. Ha{\ss}ler, Y. Matiyasevich and D. Schleicher,
\emph{Diophantine equations over \(\Z\): universal bounds and parallel
formalization}, arXiv:2506.20909v3, 22 October 2026.
Definitions 2.1, 4.1 and 5.3; Lemmas 5.1 and 6.1; Theorems 3 and A.
\url{https://arxiv.org/abs/2506.20909v3}.

\bibitem{Sun}
Z.-W. Sun, \emph{Further results on Hilbert's Tenth Problem},
Science China Mathematics \textbf{64} (2021), 281--306.
\url{https://arxiv.org/abs/1704.03504}.

\bibitem{Jones}
J. P. Jones, \emph{Universal Diophantine equation},
Journal of Symbolic Logic \textbf{47} (1982), 549--571.

\bibitem{Encoder}
\(\Z\Q\) Diophantine Project,
\emph{One rational coordinate for a finite family of algebraic witnesses},
23 November 2026.
\nolinkurl{undecidability/papers/joint_finite_algebra_encoder.tex}.

\bibitem{Sharing}
\(\Z\Q\) Diophantine Project,
\emph{A common numerator point for simultaneous integer certificates
and rational undecidability in 83, 101, and 53 variables},
23 November 2026.
\nolinkurl{undecidability/papers/common_numerator_transfers.tex};
source and degree certificate:
\nolinkurl{undecidability/certificates/common_numerator_transfers.json}.
\end{thebibliography}
\end{document}
